% theme: digestion_and_absorption
\MajorQuestion{消化と吸収}
\SetRowSpaceQanda

\Instruction{次の問いに答えなさい。}
\QQ{炭水化物、タンパク質、脂肪のように、炭素を含む養分を何というか。}{有機物}
\QQ{小腸の内側に多数ある、小さな突起を\NamedBlank{villus}という。}{柔毛}
\QQ{食物が通る、口から肛門まで続く1本の長い管を何というか。}{消化管}
\QQ{消化液に含まれ、食物の養分を分解するはたらきをもつ物質を\NamedBlank{enzyme}という。}{消化酵素}
\QQ{胆汁をつくる器官を何というか。}{肝臓}

\Instruction{\strut}
\QQ{食物の養分を、体内に取り入れやすい物質に変化させることを\NamedBlank{digestion}という。}{消化}
\QQ{ブドウ糖とアミノ酸は、\RefBlank{villus}の中の何に入るか。}{毛細血管}
\QQ{ビタミンや無機塩類のように、炭素を含まない養分を何というか。}{無機物}
\QQ{だ液に含まれ、デンプンを分解する\RefBlank{enzyme}を何というか。}{アミラーゼ}
\QQ{胆汁をたくわえる器官を何というか。}{胆のう}

\Instruction{\strut}
\QQ{\RefBlank{villus}が多数あることで、小腸の何が大きくなるか。}{表面積}
\QQ{\RefBlank{digestion}された養分を体内に取り入れることを何というか。}{吸収}
\QQ{消化管を、口から肛門まで順に6つ答えなさい。}{口、食道、胃、小腸、大腸、肛門}
\QQ{胃液に含まれ、タンパク質を分解する\RefBlank{enzyme}を何というか。}{ペプシン}
\QQ{デンプンは、最終的に何という養分に分解されるか。}{ブドウ糖}

\Instruction{\strut}
\QQ{脂肪酸とモノグリセリドは、\RefBlank{villus}の中の何に入るか。}{リンパ管}
\QQ{だ液をつくる器官を何というか。}{だ液腺}
\QQ{肝臓でつくられ、脂肪の消化を助ける消化液を何というか。}{胆汁}
\QQ{\RefBlank{enzyme}は、約何\,${}^\circ\mathrm{C}$でよくはたらくか。}{約40\,${}^\circ\mathrm{C}$}
\QQ{水分は、おもにどこで吸収され、一部はどこで吸収されるか。}{おもに小腸、一部は大腸}
